% MSc Dissertation - A Visual Trace Map of Edinburgh Place Names in Literature
% Merritt Wang | Supervisor: Uta Hinrichs | University of Edinburgh | 2026

\documentclass[logo,msc]{infthesis}
\usepackage{msccheck}
\usepackage{graphicx}
\usepackage{subcaption}
\usepackage{float}
\usepackage{longtable}
\usepackage{tabularx}
\usepackage{wasysym}
\usepackage{tikz}
\usepackage{pdfpages}
\usepackage{microtype}
\usepackage[hidelinks]{hyperref}
\usepackage{booktabs}
\usepackage{xcolor}

\begin{document}
\begin{preliminary}

\title{A Visual Trace Map of Edinburgh Place Names in Literature}
\author{Merritt Wang}
\date{\today}

\abstract{
% TODO (~200-300 words)
% Suggested structure:
%   Sentence 1-2: State the problem — existing tools visualise literary
%     place names as geographic data, but literary meaning comes from
%     narrative context, not coordinates.
%   Sentence 3-4: Introduce the data — LitLong Edinburgh database:
%     548 works (1687-2015), 2,135 place names, 48,056 mention records,
%     345 authors.
%   Sentence 5-6: State the two research directions —
%     (1) author spatial fingerprints and (2) narrative topology networks.
%   Sentence 7-8: Describe the method — interactive D3.js visualisations
%     driven by real data; user study with two participant groups.
%   Sentence 9-10: State the key findings and contribution.
}

\maketitle

\newenvironment{ethics}
   {\begin{frontenv}{Research Ethics Approval}{\LARGE}}
   {\end{frontenv}\newpage}

\begin{ethics}
This project obtained approval from the Informatics Research Ethics committee.\\
Ethics application number: 446635\\
Date when approval was obtained: % TODO: fill in date\\
The participants' information sheet and a consent form are included in the appendix.\\
\standarddeclaration
\end{ethics}

\begin{acknowledgements}
% TODO
\end{acknowledgements}

\tableofcontents
\listoffigures

\end{preliminary}


% ============================================================
% CHAPTER 1: INTRODUCTION (~1,500 words)
% ============================================================
\chapter{Introduction}

\section{Motivation}
% TODO (~300 words)
%
% Key points to cover:
% - Edinburgh is one of the most literarily documented cities in the world.
%   The LitLong / Palimpsest project (University of Edinburgh) extracted
%   place-name mentions from 548 literary works spanning 1687-2015,
%   producing 48,056 geo-referenced records.
%
% - Existing tools (including the original LitLong interface) visualise
%   this data on geographic maps. This assumes that a place name's meaning
%   is determined by its location.
%
% - In literature, meaning comes from narrative context:
%   WHEN in the narrative a place appears (early/late in the story),
%   WHAT other places surround it (co-occurrence),
%   WHO chose it (authorial style),
%   HOW frequently it recurs.
%
% - A geographic map cannot show any of this. It can only show WHERE.
%
% - Motivating example to consider using:
%   Irvine Welsh's Trainspotting clusters almost entirely around Leith
%   (41% of all mentions in the Leith & North sector), while
%   Walter Scott clusters around Old Town (33.5%).
%   On a geographic map, both authors' dots appear in the same city.
%   Only a purpose-built visualisation can reveal this divergence.

\section{Problem Statement}
% TODO (~200 words)
%
% Core question: Can place names in literary text be visualised as an
% independent narrative structure, rather than simply as a rendering of
% geographic information?
%
% Two sub-problems:
% (1) Can an author's spatial attention be represented as a distinctive
%     visual "fingerprint" that allows different authors to be distinguished?
% (2) Do places that co-occur in literary works form a narrative topology
%     that diverges meaningfully from their geographic topology?
%
% What existing tools do NOT do:
% - They do not separate literary space from geographic space.
% - They do not visualise authorial spatial style.
% - They do not show which places are narratively proximate vs
%   geographically proximate.

\section{Research Questions}
% TODO (~200 words)
%
% RQ1 (Direction 1 — Author Spatial Fingerprint):
%   Does each author's distribution of place-name mentions across
%   Edinburgh's city sectors form a visually distinctive spatial signature
%   that allows users to identify authors from their fingerprint alone?
%
% RQ2 (Direction 2 — Narrative Topology):
%   Does the co-occurrence of place names within literary works reveal a
%   narrative topology that diverges from the geographic topology of
%   Edinburgh?
%
% Success criterion for RQ1 (defined by supervisor):
%   Users should be able to match a fingerprint to an author on a
%   one-to-one basis — i.e. different authors must produce visually
%   distinct shapes (one-to-one, not many-to-one).
%
% Success criterion for RQ2:
%   The visualisation should reveal cases where narrative proximity
%   contradicts geographic proximity — e.g. Leith and Princes Street
%   are the strongest co-occurring pair (weight=28) despite being
%   geographically separated.

\section{Contributions}
% TODO (~200 words)
%
% List of contributions:
% 1. A novel interactive visualisation system for literary place-name data
%    that separates narrative structure from geographic information.
%
% 2. Two complementary visualisation directions implemented in D3.js:
%    author spatial fingerprints and narrative topology networks.
%
% 3. An empirical finding with implications for literary geography:
%    among 5 representative Edinburgh authors (McCall Smith, Welsh,
%    Lockhart, Stevenson, Scott), ZERO place names are shared by 3 or
%    more authors. Only 18 place names are shared by exactly 2 authors.
%    This confirms that each author constructs an almost entirely private
%    imaginary Edinburgh.
%
% 4. A validated method for extracting narrative order from the LitLong
%    database using the start_word field, which is globally unique and
%    monotonically increasing within each document.
%
% 5. A user evaluation comparing domain expert and general public responses
%    to multiple visualisation designs, informing the final design selection.

\section{Dissertation Overview}
% TODO (~100 words)
% Brief roadmap: what each chapter covers.


% ============================================================
% CHAPTER 2: LITERATURE REVIEW (~3,000 words)
% ============================================================
\chapter{Literature Review}

\section{Literary Geography and Geocriticism}
% TODO (~500 words)
%
% Key references to cover:
% - Westphal (2011): Geocriticism — place as a cultural and narrative
%   construct, not a fixed coordinate.
% - Tally (2013): Spatiality — the "spatial turn" in literary studies.
% - Franco Moretti (2005): Graphs, Maps, Trees — quantitative approaches
%   to literary geography; his literary maps of European novels.
%   Important: Moretti himself acknowledges that his maps reveal the
%   ABSENCE of places as much as their presence (useful for your
%   "literary silence" angle if you develop Direction 4 later).
%
% Key argument to make:
%   Literary geography has established that place in fiction is not
%   simply a geographic referent but a narrative and cultural construction.
%   However, most digital literary geography tools have not caught up with
%   this theoretical insight — they still plot place names on coordinates.

\section{The LitLong Edinburgh Project}
% TODO (~400 words)
%
% Key points:
% - Anderson et al. (2016): Palimpsest / LitLong project description.
%   The NLP pipeline used to extract place names from literary texts.
% - Database structure: 620 documents, 2,135 place names, 50,248 mention
%   records, 424 authors, 29 literary genres.
%   (Note: the SQL dump is more complete than the CSV snapshot —
%   620 docs vs 548 in CSV, 50,248 mentions vs 48,056.)
% - The api_sentence table contains full source text for each mention —
%   this enables the "hover to read original sentence" feature.
% - The start_word field encodes narrative position and can be used
%   to reconstruct reading order within each document.
% - Prior visualisation: the original LitLong web interface used a
%   geographic map — describe its limitations.

\section{Visualising Place in Literary Corpora}
% TODO (~500 words)
%
% Key references:
% - Stange & Dork (2015): "Unfolding the City" / Novel City Maps.
%   CRITICAL reference for Direction 1. They created street-map-style
%   fingerprints of literary place mentions. Your work pushes further:
%   you DROP geographic coordinates entirely and use abstract sectors.
%
% - Google Literary Atlas / other mapping projects — what they do and
%   why they are still fundamentally geographic.
%
% - Cooper et al. on literary GIS — the tension between geographic
%   precision and literary interpretation.
%
% Key argument: existing visualisation work either uses geographic maps
% (losing narrative structure) or is fully abstract (losing geographic
% intuition). This project aims for a middle ground.

\section{Information Visualisation for Humanities Data}
% TODO (~500 words)
%
% Key references:
% - Shneiderman (1996): "Overview first, zoom and filter, details on demand"
%   — the information seeking mantra that underpins your interactive design.
% - Tufte (1983): small multiples as a principle — directly relevant to
%   your Small Multiples visualisation.
% - Drucker (2011): "Humanities Approaches to Graphical Display" —
%   argues that humanistic data is "capta" (interpreted) not "data"
%   (given), and that visualisations should reflect uncertainty and
%   interpretation. Relevant to your design philosophy.
% - Hinrichs & Carpendale (2011): "Gestures in the Wild" — your supervisor's
%   own work on embodied interaction with visualisation; worth citing.

\section{Graph and Network Visualisation}
% TODO (~500 words)
%
% Key references:
% - Fruchterman & Reingold (1991): force-directed graph layout algorithm
%   used in your network visualisation.
% - Nollenburg & Wolff (2006): metro-map layout algorithm —
%   CITE THIS in relation to why you used manual layout instead of
%   algorithmic: "Schematic metro-map layout is a known hard problem
%   in graph drawing (Nollenburg & Wolff, 2006)."
% - Beck et al. (2017): survey of set visualisation techniques.
% - London Tube Map (Beck, 1933) as a design precedent for replacing
%   geographic topology with functional/narrative topology.
%   Key point: Beck's map was initially REJECTED because it "didn't look
%   like a map" — but it succeeded because it served the traveller's
%   actual need (which line, which direction) rather than geographic
%   accuracy. Your narrative topology serves the reader's need
%   (which places are narratively connected) rather than geographic fidelity.

\section{Gap and Contribution}
% TODO (~300 words)
%
% Synthesise the gap:
% - Literary geography theory says place is narrative, not geographic.
% - But existing digital tools still use geographic maps.
% - Stange & Dork came closest but kept geographic coordinates.
% - No existing tool visualises place co-occurrence as narrative topology.
% - No existing tool allows direct comparison of author spatial styles
%   using abstract (non-geographic) fingerprints.
% This project fills that gap.


% ============================================================
% CHAPTER 3: DESIGN (~2,000 words)
% ============================================================
\chapter{Design}

\section{Design Rationale}
% TODO (~300 words)
%
% Why interactive:
% - Static maps cannot support the "overview + detail" interaction
%   that humanities scholars need to explore a 548-work corpus.
% - Interactivity allows users to filter by author, adjust thresholds,
%   and access source text on demand.
%
% Why D3.js:
% - Proposed in the project brief.
% - Web-native: no installation required, results shareable via URL.
% - R's interactive options (Shiny, Plotly) require a server and are
%   not web-native. Supervisor noted uncertainty about R's interactive
%   capability in Week 4 meeting.
% - D3.js is the standard tool in digital humanities visualisation.
%
% Design philosophy:
% - Narrative structure over geographic accuracy.
% - The visualisations deliberately move away from coordinates.
%   Where coordinates are used (Small Multiples), they serve as a
%   loose spatial anchor, not a precise map.

\section{Data Exploration and Key Findings}
% TODO (~400 words)
%
% IMPORTANT — report these findings with specific numbers:
%
% Finding 1 — Author spatial languages are highly individual:
%   Among the 5 test authors, ZERO place names are shared by 3 or more.
%   Only 18 place names are shared by exactly 2 authors.
%   This means individual place names cannot serve as radar chart axes
%   (most axes would be zero for most authors).
%   Decision: use city sectors as axes instead (supervisor's suggestion).
%
% Finding 2 — Sector distributions validate the approach:
%   - McCall Smith: New Town 60.7% (Scotland Street series set there)
%   - Welsh: Leith & North 41.1% (Trainspotting centred on Leith)
%   - Scott: Old Town 33.5% (historical fiction in Old Town)
%   - Lockhart: New Town 44.3% (19th century literary circle)
%   - Stevenson: New Town 35.2% + Leith 19.8% (spans both areas)
%   Data-driven result matches established literary-geographic knowledge.
%   This cross-validation strengthens the method's credibility.
%
% Finding 3 — Co-occurrence granularity:
%   Sentence-level: 0 co-occurring pairs (too fine-grained).
%   Page-level: up to 92 co-occurring places per page (too coarse, hairball).
%   Document-level: 403 pairs with weight >= 2 (selected as appropriate).
%
% Finding 4 — Narrative topology diverges from geographic topology:
%   Strongest co-occurring pair: Leith & Princes Street (weight=28).
%   These are geographically separated (Leith is a port district north
%   of the city centre; Princes Street is the main shopping street).
%   Yet they appear together in 28 different literary works.
%   This is a concrete example of narrative proximity != geographic proximity.
%
% Finding 5 — Walter Scott's broad geographic range:
%   Scott's data includes locations far outside Edinburgh city centre
%   (Forth: lat 56.05, Haddington: lon -2.80, Linlithgow: lon -3.58).
%   This caused coordinate range compression in Small Multiples.
%   Decision: fix coordinate range to Edinburgh core area
%   (lat 55.88-56.01, lon -3.30 to -3.05).

\section{Direction 1: Author Spatial Fingerprint}
% TODO (~500 words)
%
% Design candidates considered (with rationale for each):
%
% (a) Radar Chart — SELECTED as primary design
%     Axes = 6 city sectors (not individual place names — see Finding 1).
%     Shape = author's fingerprint.
%     Why city sectors: supervisor's hand-drawn sketch explicitly noted
%     "city sectors"; data confirmed this is the only viable axis strategy.
%     Sector definitions based on lat/lon coordinate ranges derived from
%     api_location data.
%
% (b) Bar-code Fingerprint — retained as secondary design
%     Each place = one vertical bar; height = mention frequency (log scale).
%     Why log scale: Welsh's Leith mention count (650) is an extreme outlier.
%     On linear scale, all other bars compressed flat. Log scale is standard
%     treatment for long-tail distributions.
%
% (c) Small Multiples — retained as tertiary design
%     5 panels side-by-side; real-world coordinates (supervisor's sketch
%     noted "Real-world"); abstract city map overlay for spatial context.
%     Finding: single-panel readability is lower than radar chart;
%     value is comparative across panels.
%
% (d) Colour Map Silhouette — explored but deprioritised
%     City boundary silhouette with heat overlay.
%     Issue: if all authors concentrate in city centre, discriminability
%     is low. Retained as a future direction.

\section{Direction 2: Narrative Topology Network}
% TODO (~500 words)
%
% Design candidates considered:
%
% (a) Force-directed Network — SELECTED as primary design
%     Nodes = places (size = mention frequency).
%     Edges = document-level co-occurrence (weight = number of works).
%     Layout: D3 force simulation.
%     Node position carries NO geographic meaning — proximity in graph
%     = narrative proximity only.
%     Key debugging decisions:
%       v1: weight * 0.3 — no visible width difference.
%       v2: sqrt scaling — marginal improvement.
%       v3: weight >= 8 threshold + (w/max)^2 * 12 scaling + k=3.5.
%       v3 resolved both density and scaling issues.
%
% (b) Linear Connection Diagram — retained as secondary design
%     Places on a linear axis; Bezier arcs for co-occurrence.
%     REASON for linear over chord (circular):
%       Supervisor feedback: circular form implies cyclical structure,
%       which is inappropriate for literary narratives (which are linear).
%       "The book has a beginning and an end."
%     Arc height proportional to node distance; weight >= 6 threshold.
%
% (c) Metro-map Topology — retained as conceptual/illustrative design
%     8 lines, 50+ stations on manual grid layout.
%     Echoes the London Underground map (Beck, 1933) cited in project brief.
%     Why manual layout: algorithmic schematic layout is a known hard
%     problem (Nollenburg & Wolff, 2006); deferred for future work.
%
% (d) Adjacency Matrix Heatmap — explored but deprioritised
%     Most precise representation but unfriendly for general audiences.
%     May be appropriate for domain-expert-only study condition.
%
% (e) Chord Diagram — replaced by linear diagram
%     See supervisor feedback above.

\section{Sketch Evaluation and Design Selection}
% TODO (~300 words)
%
% Describe the user study approach (detail in Chapter 5):
% - 6 candidate designs presented to participants without over-explanation.
% - Two tasks: (1) author identification task (fingerprint direction),
%   (2) narrative relationship identification task (topology direction).
% - Vote on most intuitive design.
% - Two participant groups: domain experts vs general public.
% - Hypothesis: experts may prefer precise designs (adjacency matrix,
%   bar-code); general public may prefer intuitive designs (radar chart,
%   metro map).
% - Final design selection based on task accuracy rates + participant votes.


% ============================================================
% CHAPTER 4: IMPLEMENTATION (~3,500 words)
% ============================================================
\chapter{Implementation}

\section{Data Pipeline}
% TODO (~500 words)
%
% Step-by-step pipeline to describe:
%
% 1. PostgreSQL setup:
%    - PostgreSQL 16 installed via Homebrew on macOS.
%    - Full LitLong dump (litlong_2022-06-21.sql, 145MB, 2.3M lines)
%      imported directly. This is the complete official database, more
%      recent and complete than the CSV snapshots.
%    - Issue: api_location table depends on PostGIS geometry types
%      (geom, poly fields). PostGIS not installed. Solution: imported
%      api_location from CSV snapshot, retaining only id, text, lat, lon,
%      gazref, ptype. PostGIS fields (geom, poly) not needed for current
%      visualisation directions.
%
% 2. Mention order integration:
%    - start_word field format: "w81489" (global word index).
%    - Validated: globally unique, strictly monotonically increasing
%      within each document (tested on documents 1, 106, 250).
%    - New table mention_order created with fields:
%        id, document_id, location_id, word_num (integer),
%        mention_order (rank within document),
%        position_pct (normalised 0-1 position).
%    - Formula: position_pct = (word_num - min) / (max - min)
%      per document_id group.
%    - Validation: cross-checked against source text in api_sentence table.
%      api_sentence contains 50,248 records (one per locationmention).
%
% 3. Six-table JOIN query validated:
%    author + title + place + mention_order + position_pct + source sentence
%    all retrievable in a single SQL query.
%
% 4. Key database statistics after full import:
%    api_document: 620 records
%    api_locationmention: 50,248 records
%    api_sentence: 50,248 records (full source text)
%    api_location: 2,135 records
%    api_genre: 29 literary genres
%    api_author: 424 authors
%    api_posmention: 2,085,834 POS annotation records (future work)

\section{Direction 1: Author Spatial Fingerprint — Implementation}
% TODO (~800 words)
%
% Radar Chart:
% - D3.lineRadial with curveLinearClosed.
% - Y-axis scaled to data maximum (0.65) not 1.0:
%   McCall Smith's New Town value (60.7%) is the highest;
%   fixing to 100% left 35% of the chart empty.
% - Hover targets: invisible circles r=14px (not r=6px) for comfortable
%   interaction without pixel-precise cursor placement.
% - Legend: click to show/hide individual authors.
% - Data embedded directly in HTML (no external CSV dependency):
%   reason: browsers block local file reads under file:// protocol.
%
% Bar-code Fingerprint:
% - Per-author SVG rows; bars = top 30 place names by combined frequency.
% - Log scale: np.log1p(value * 1000). Resolves Welsh/Leith outlier problem.
% - Hover: place name + percentage value.
%
% Small Multiples:
% - 5 panels in 3+2 flex layout.
% - Abstract city map SVG overlay on each panel:
%     Firth of Forth (blue water polygon)
%     City landmass (beige polygon with approximate coastline)
%     Arthur's Seat (ellipse + hill silhouette path)
%     Pentland Hills (ridge-line polygon)
%     Edinburgh Castle (rectangular tower symbol)
%     Holyrood (diamond symbol)
%     Water of Leith (curved blue path)
%     Neighbourhood labels (Old Town, New Town, Leith, Morningside, Stockbridge)
% - Coordinate range fixed to Edinburgh core:
%     lat 55.88-56.01, lon -3.30 to -3.05.
%     Points outside range filtered out.
%     Reason: Scott's broad geographic range compressed city centre.
% - Per-panel zoom (D3 zoom, scale 1-10x) on a zoomG group,
%   allowing independent inspection of dense clusters.

\section{Direction 2: Narrative Topology — Implementation}
% TODO (~800 words)
%
% Force-directed Network:
% - D3 force simulation:
%     forceLink: distance inversely proportional to weight (stronger
%       co-occurrence = shorter distance).
%     forceManyBody: strength -400 (repulsion).
%     forceCollide: radius = nodeRadius + 18 (prevents overlap).
%     forceCenter: pulls graph to canvas centre.
% - Node radius: log1p(total_mentions) * 2.8.
% - Edge width: (weight/maxWeight)^2 * 12 (squared normalisation).
% - Weight threshold slider (range 8-28): simulation re-run on change.
% - Labels rendered BELOW nodes (y = node.y + radius + 12):
%     Initial bug: labels rendered inside nodes with fill=white,
%     invisible against dark background.
% - Hover: top 5 co-occurring places for selected node.
% - Drag nodes: D3 drag behaviour with fx/fy pinning.
% - Zoom/pan: D3 zoom on outer SVG, transform applied to inner group.
%
% Linear Connection Diagram:
% - Places ordered left-to-right by total mention frequency.
% - Bezier arcs: M x1 y0 Q midX (y0-height) x2 y0.
%   Arc height = horizontal distance * 0.42 (proportional).
% - Stroke width: (weight/maxWeight)^1.5 * 6.
% - Weight threshold slider (range 6-28).
% - Hover arcs: source, target, weight.
% - Hover nodes: total mention count.
%
% Metro-map:
% - 8 lines, 50+ stations on manual grid (unit = 58px).
% - Interchange stations (served by 2+ lines) marked with double-circle.
% - Zoom/pan via D3 zoom behaviour.
% - Hover: station name, lines served, interchange indicator.
% - Why manual layout: algorithmic metro-map layout (Nollenburg & Wolff,
%   2006) is a known NP-hard graph drawing problem; deferred.

\section{Combined Interface}
% TODO (~600 words)
%
% FUTURE WORK — to be implemented Week 6-7.
%
% Planned architecture:
% - Single-page HTML application.
% - Top navigation: toggle between "Fingerprints" and "Topology" views.
% - Left control panel:
%     Author selector (multi-select dropdown, all 424 authors available,
%     default: 5 test authors).
%     Visualisation type selector (within each direction).
%     Parameter sliders (co-occurrence threshold, etc.).
% - Main view: selected visualisation rendered in centre.
% - Detail panel (right or bottom):
%     On hover/click of any place name: display source sentences from
%     api_sentence table — this is the "hover to read original text" feature.
%     This requires a backend query or pre-loading sentence data.
% - Cross-view linking:
%     Clicking an author in the fingerprint view highlights their places
%     in the topology view.
%
% Scale at combined interface stage:
% - Expand author pool from 5 to 20-30 representative authors.
% - No need to include all 424 — select for diversity in time period,
%   genre, and geographic focus.
% - Network graph: expand to top 50 high-frequency place names.

\section{Technical Decisions and Constraints}
% TODO (~300 words)
%
% Data embedding:
% - All D3 HTML files embed data as JavaScript arrays.
% - Reason: d3.csv() fails under file:// protocol (browser security).
% - Trade-off: larger HTML files; data cannot be updated without
%   re-embedding. Acceptable for dissertation prototype.
%
% D3.js version:
% - D3 v7 loaded from CDN (d3js.org/d3.v7.min.js).
% - Confirmed accessible (HTTP 200) from the development environment.
%
% PostgreSQL vs CSV:
% - PostgreSQL used as primary data store.
% - CSV snapshots retained as backup and for reproducibility.
% - mention_order table exported as SQL for archiving.
%
% PostGIS:
% - Not installed (geom/poly fields excluded from api_location).
% - Impact: precise polygon boundaries for place areas not available.
% - Mitigation: lat/lon centroids are sufficient for all current
%   visualisation directions.


% ============================================================
% CHAPTER 5: EVALUATION (~2,500 words)
% ============================================================
\chapter{Evaluation}

\section{Study Design}
% TODO (~400 words)
%
% Key design decisions (decided during Week 4 meeting preparation):
%
% Two participant groups:
%   Group A — Domain experts (geography / literature / digital humanities)
%   Group B — General public
%   Rationale: hypothesis that design preferences differ between groups.
%   Finding from this distinction will be reported as a contribution.
%
% Online format:
%   Reasons: (1) broader recruitment beyond university;
%             (2) larger sample size;
%             (3) saves participant and researcher time.
%   Tool: consider Qualtrics (university licence available) or Google Forms.
%
% Introduction protocol (IMPORTANT):
%   Participants will be shown visualisations WITHOUT detailed explanation
%   of what each chart means. Only the project purpose is explained.
%   Reason: avoids "priming" participants; tests whether designs are
%   self-explanatory. This was Merritt's own decision, confirmed as
%   methodologically sound.
%
% Two task types:
%   Task 1 (fingerprint direction): Author identification task.
%     Show participant a fingerprint (radar chart / bar-code / small multiples).
%     Ask: "Which of these authors do you think this pattern belongs to?"
%     Metric: accuracy rate (proportion of correct identifications).
%   Task 2 (topology direction): Narrative relationship task.
%     Show participant a network / linear diagram.
%     Ask: "Which two places seem most narratively connected in this diagram?"
%     Also ask: "Are these places geographically close or far apart?"
%     Metric: accuracy + whether user perceives narrative/geographic divergence.
%
% Vote:
%   After seeing all designs: "Which design felt most intuitive?"
%   Separate votes for fingerprint direction and topology direction.

\section{Participants}
% TODO (~200 words)
% - Recruitment method (social media, university mailing lists, etc.)
% - Target sample size (suggest: minimum 20 per group, 40 total)
% - Demographics collected (age range, background, prior experience
%   with data visualisation, familiarity with Edinburgh)
% - Inclusion/exclusion criteria

\section{Study Procedure}
% TODO (~400 words)
% - Consent process (online consent form, based on ethics approval 446635)
% - Introduction: project purpose explained, designs NOT explained
% - Task sequence
% - Estimated duration (~15-20 minutes)
% - Data collected: task accuracy, response time (if measurable),
%   design vote, open-ended comments

\section{Results}
% TODO (~800 words)
%
% Suggested structure:
% 3.1 Task 1 results: author identification accuracy per design
%     (which fingerprint design allowed most accurate identification?)
% 3.2 Task 2 results: narrative relationship identification accuracy
%     (did participants notice narrative/geographic divergence?)
% 3.3 Design votes: most intuitive per direction
% 3.4 Expert vs general public comparison
% 3.5 Qualitative themes from open-ended comments

\section{Discussion of Results}
% TODO (~400 words)
%
% Questions to address:
% - Which design succeeded best at the author identification task?
%   (Prediction: radar chart, because shape differences are salient)
% - Did participants spontaneously notice that Leith and Princes Street
%   are strongly co-occurring but geographically separated?
%   (This would validate the core thesis of RQ2)
% - Were there significant differences between expert and general public?
% - Were there any surprising results?


% ============================================================
% CHAPTER 6: DISCUSSION (~2,000 words)
% ============================================================
\chapter{Discussion}

\section{Design Improvements after User Study}
% TODO (~600 words)
%
% Address: what changed in the final designs based on user feedback?
% - Which visualisations were promoted to the combined interface?
% - Which were demoted to appendix / supplementary material?
% - What specific visual elements were changed?

\section{Narrative Topology vs Geographic Topology}
% TODO (~500 words)
%
% Develop the key finding:
% - Leith and Princes Street: co-occurrence weight 28, strongest pair.
%   Geographically: Leith is a port district ~2.5km north of Princes Street.
%   Narratively: they appear together in 28 different literary works.
% - What does this mean? Authors move characters between Leith and
%   Princes Street as a recurrent narrative pattern. The "journey" between
%   these two places — between the working-class port and the bourgeois
%   shopping street — is a structuring device in Edinburgh fiction.
% - Connect to Welsh: Trainspotting explicitly dramatises this journey
%   as a journey between social worlds, not just geographic locations.
% - This is the kind of insight that a geographic map cannot reveal.
%   The narrative topology visualisation makes it visible.

\section{Author Spatial Languages}
% TODO (~500 words)
%
% Develop the key finding:
% - Zero shared places among 5 authors at threshold >= 3.
%   Only 18 shared at threshold = 2.
% - This quantifies what literary scholars have described qualitatively:
%   each author has a distinctive imaginary Edinburgh.
% - McCall Smith's Edinburgh is bounded by the New Town triangle
%   (Scotland Street, Dundas Street, Drummond Place); Welsh's Edinburgh
%   is bounded by Leith.
% - The fingerprint visualisation makes this visible and comparable
%   across authors simultaneously.
% - Future work: expand to all 424 authors in the database to produce
%   a complete "atlas of Edinburgh literary imaginations."

\section{Limitations and Open Questions}
% TODO (~400 words)
%
% Honest limitations to address:
%
% 1. Manual metro-map layout:
%    Algorithmic schematic layout deferred (NP-hard problem).
%    Current layout is hand-crafted, not generalisable to arbitrary data.
%
% 2. PostGIS not available:
%    Precise polygon boundaries for places excluded.
%    Lat/lon centroids used instead. Some places (e.g. "Old Town")
%    have centroids that may not accurately represent the full area.
%
% 3. 5 test authors is a small sample:
%    Sufficient for proof of concept; insufficient for generalisation.
%    Full deployment with 424 authors would require UI scaling work
%    (dropdown selectors, pagination, clustering).
%
% 4. Reader-sketch direction not fully explored:
%    Supervisor mentioned "Reader-sketch?" with a question mark in notes.
%    This direction (visualising reading experience / narrative time)
%    was not fully developed in this project.
%    The mention_order / position_pct data is ready for this.
%    The Reader-plot timeline (linear axis showing place order across
%    a book) was prototyped but not included in the user study.
%
% 5. api_posmention not used:
%    2,085,834 POS annotation records available but not incorporated.
%    Could support analysis of whether a place appears in dialogue
%    vs narration (potential Direction 3 / narrative weight analysis).
%
% 6. Source text coverage:
%    directus.sql contains curated source text excerpts for some places;
%    api_sentence has full sentences for all 50,248 mentions.
%    The "hover to read original sentence" feature is planned but not
%    yet implemented in the combined interface.


% ============================================================
% CHAPTER 7: CONCLUSION (~700 words)
% ============================================================
\chapter{Conclusion}

% TODO (~700 words)
%
% Structure:
% 1. Restate the problem and why it matters.
% 2. Summarise what was built and what was found.
% 3. Answer the two research questions explicitly:
%    RQ1: Yes/No/Partial — did the fingerprint allow author identification?
%    RQ2: Yes/No/Partial — did the topology diverge from geography?
% 4. State the three main contributions.
% 5. Future work (2-3 paragraphs):
%    a. Expand to all 424 authors (combined interface with dropdown).
%    b. Implement Reader-plot timeline using position_pct data.
%    c. Investigate Direction 3: narrative weight using api_posmention
%       (does a place appear in dialogue vs narration? in climax vs
%       background description?).
%    d. Explore "literary silences" — which Edinburgh places NEVER appear
%       in literature? (Direction 4, mentioned in proposal).
%    e. Algorithmic metro-map layout as a graph drawing research problem.


% ============================================================
% BIBLIOGRAPHY
% ============================================================
\bibliographystyle{plain}
\bibliography{mybibfile}
%
% Key references to include (start collecting now):
% - Anderson et al. (2016): LitLong / Palimpsest project paper
% - Westphal (2011): Geocriticism
% - Tally (2013): Spatiality
% - Moretti (2005): Graphs, Maps, Trees
% - Stange & Dork (2015): Novel City Maps (CRITICAL)
% - Nollenburg & Wolff (2006): Metro-map layout algorithm
% - Shneiderman (1996): Information seeking mantra
% - Tufte (1983): Envisioning Information (small multiples)
% - Drucker (2011): Humanities approaches to graphical display
% - Fruchterman & Reingold (1991): Force-directed layout
% - Hinrichs & Carpendale (2011): supervisor's own work
% - Beck (1933): London Underground map (design precedent)


% ============================================================
% APPENDICES
% ============================================================
\appendix

\chapter{Study Material}

\section{Participant Information Sheet}
% TODO: insert PIS PDF
% \includepdf[pages=-]{pis.pdf}

\section{Consent Form}
% TODO: insert Consent Form PDF
% 
\includepdf[pages=-]{consent.pdf}

\section{Task Questions}
% TODO

\section{Survey / Interview Questions}
% TODO

\chapter{Data and Implementation}

\section{Database Schema}
% TODO: reference or include litlong_schema.sql
% Key tables: api_document, api_locationmention, api_sentence,
%             api_location, api_genre, api_author,
%             api_document_author, mention_order (custom)

\section{City Sector Definitions}
% TODO: table of lat/lon boundary ranges for 6 sectors:
%   Leith & North:       lat > 55.97
%   New Town:            lat >= 55.952, lon in [-3.22, -3.175]
%   Old Town:            lat < 55.952, lon in [-3.20, -3.165]
%   East & Arthur's Seat: lon > -3.165, lat < 55.96
%   South Edinburgh:     lat < 55.945
%   Outer Edinburgh:     lat < 55.92 or lon < -3.28

\section{Author Selection}
% TODO: table of 5 test authors:
%   McCall Smith: 7,320 mentions, 17 books
%   Welsh:        2,634 mentions, 8 books
%   Lockhart:     3,046 mentions, 6 books
%   Scott:        2,796 mentions, 12 books
%   Stevenson:    1,778 mentions, 16 books

\section{Co-occurrence Analysis}
% TODO: report the 403 document-level co-occurring pairs (weight >= 2)
% Top 10 pairs table; granularity decision table (sentence/page/document)

\end{document}